
\subsubsection{2.1 Spurious Correlations and Group Robustness}

Spurious correlations --- statistical associations between features and labels that do not reflect causal structure --- cause ERM-trained models to fail systematically on minority groups \citep{sagawa2019distributionally}. The Waterbirds benchmark \citep{sagawa2019distributionally} isolates this phenomenon: bird labels are spuriously correlated with backgrounds, creating four groups (2 class $\times$ 2 background) with dramatically different ERM performance. Distributionally robust optimization (DRO) \citep{sagawa2019distributionally} directly targets worst-group accuracy but requires group annotations at training time.

\subsubsection{2.2 Annotation-Required Methods: DFR}

The most effective post-hoc approach is deep feature reweighting (DFR) \citep{kirichenko2022last}: retrain only the last layer on a balanced held-out validation set, keeping ERM-pretrained features frozen. DFR achieves worst-group accuracy (WGA) approximately 88--90\% on Waterbirds with a small balanced set. Hill et al. [2025] explain DFR's effectiveness as implicit group-balancing: when the retraining set overrepresents minority samples, the final classifier learns to rely on causal features rather than spurious shortcuts. DFR's key limitation is the requirement for explicit group annotations in the balanced validation set.

\subsubsection{2.3 Annotation-Free Proxies: First-Order Signals}

The annotation-free direction emerged with JTT \citep{liu2021just}: identify misclassified samples after a first ERM run, upweight them, and retrain. JTT achieves WGA approximately 80--84\% on Waterbirds without group labels. SELF \citep{labonte2023towards} refines this by combining misclassification-based proxy selection with self-training, reaching WGA approximately 82--87\%. All such methods share a fundamental constraint: they operate on first-order signals --- model output or loss --- and are insensitive to loss landscape geometry.

The present work proposes the first second-order annotation-free proxy: per-sample Hessian trace of the last-fc layer. This signal is orthogonal to loss and misclassification rates and carries information about differential curvature created by spurious feature exploitation.

\subsubsection{2.4 Loss Landscape Geometry and Hessian Analysis}

Sharpness of loss minima is connected to generalization: sharp minima (high Hessian eigenvalues) generalize poorly compared to flat minima \citep{keskar2016large}. Sharpness-aware minimization (SAM) \citep{foret2021sharpnessaware} explicitly seeks flat minima. Per-layer Hessian trace estimation via Hutchinson's method has been studied for neural network diagnostics \citep{yao2020pyhessian}. The Hutchinson estimator approximates tr(H) $\approx$ (1/K)$\Sigma v^T$ H v for K random Rademacher vectors v without materializing the full Hessian. However, no prior work extends Hutchinson estimation to the per-sample level to discriminate minority from majority group membership.

\subsubsection{2.5 Theoretical Connections}

LaBonte et al. [2024] show that minority group covariance matrices have larger spectral norm than majority groups within the same class --- a group-level prediction of systematic representation-norm inequality ($\|x_{i}\|^{2})$. This finding supports the hypothesis that per-sample Hessian traces would be elevated for minority samples, since Tr($H_i^{fc}$) ∝ $\|x_{i}\|^2$ p\_i($1-p_{i})$ for a linear cross-entropy head. LaBonte and Muthukumar [2026] prove that SGD on XOR spurious models learns the spurious feature first (Phase I) before the core feature (Phase II), creating a transient window where minority samples remain near the decision boundary.

This work extends LaBonte et al. [2024] from group-level covariance to per-sample trajectory analysis, and tests the Phase I/II prediction of LaBonte and Muthukumar [2026] on real image data with ResNet-50.
